\documentclass[10pt]{article}
\usepackage[margin=1in]{geometry}
\usepackage{amsmath, amssymb, amsthm}
\usepackage{enumitem}
\usepackage{hyperref}
\usepackage{booktabs}

\newtheorem*{claim}{Claim}
\newcommand{\bigO}{O}
\newcommand{\RMQ}{\operatorname{RMQ}}
\newcommand{\LCA}{\operatorname{LCA}}
\newcommand{\pred}{\operatorname{pred}}
\newcommand{\succs}{\operatorname{succ}}
\newcommand{\lcp}{\operatorname{LCP}}
\newcommand{\SA}{\operatorname{SA}}

\title{EDA Final: Ejercicios y Soluciones Propias\\
\large Basado en MIT 6.851 Spring 2021 y slides locales}
\author{}
\date{}

\begin{document}
\maketitle

\section*{Como usar esta hoja}
Cada ejercicio esta pensado para responderse a papel en una pagina: idea,
pseudocodigo, prueba corta y complejidad. Los problem sets oficiales de MIT
estan archivados localmente en \texttt{../originales/}; aqui solo se incluyen
referencias y soluciones propias/resumidas.

\section{Integer Data Structures}

\subsection*{I1. vEB: tiempo de predecessor}
Sea un universo de tamano $u = 2^w$. Una estructura van Emde Boas divide cada
clave $x$ como $x=\langle c,i\rangle$, donde $c$ es el cluster e $i$ la posicion
dentro del cluster. Explica predecessor/successor y prueba el tiempo.

\paragraph{Solucion.}
Se guarda $\min$, $\max$, un \emph{summary} de clusters no vacios y una
estructura por cluster. Para $\succs(x=\langle c,i\rangle)$:
si $x < \min$, se responde $\min$. Si el cluster $c$ existe y
$i < \max(cluster[c])$, se busca recursivamente dentro de ese cluster. Si no,
se busca en \texttt{summary} el siguiente cluster no vacio $c'$ y se responde
$\langle c', \min(cluster[c'])\rangle$. Cada llamada reduce el universo a
$\sqrt{u}$, entonces
\[
T(u)=T(\sqrt{u})+\bigO(1)=\bigO(\log\log u)=\bigO(\log w).
\]
La misma idea funciona simetricamente para predecessor.

\subsection*{I2. PS7 MIT: espacio de vEB con saving space}
Referencia: \texttt{../originales/psets/ps7.pdf}. Estudia la version que guarda
solo clusters no vacios. Da una cota de espacio y como llegar a espacio lineal.

\paragraph{Solucion propia.}
La version sparse elimina clusters vacios mediante hashing, pero una clave puede
aparecer implicitamente en varios niveles por la ruta recursiva. Una cota natural
es $\bigO(n\log\log u)$ palabras si cada elemento causa presencia en un nivel
por escala. Para espacio lineal, se usa indireccion: partir los $n$ elementos en
buckets de tamano $\Theta(\log u)$ o $\Theta(w)$ mantenidos por estructuras
simples, y construir vEB/x-fast/y-fast solo sobre representantes. El numero de
representantes baja a $\bigO(n/w)$ y cada bucket ocupa espacio proporcional a su
tamano; la suma queda $\bigO(n)$ palabras, manteniendo busquedas con el costo
del representante mas una busqueda local.

\subsection*{I3. x-fast versus y-fast}
Compara x-fast y y-fast para predecessor en word RAM.

\paragraph{Solucion.}
x-fast guarda todos los prefijos de cada clave en tablas hash por nivel. Para
buscar una clave $x$, hace binary search sobre la longitud del prefijo mas largo
presente y luego usa punteros de hojas para obtener predecessor/successor.
Tiempo $\bigO(\log w)$, espacio $\Theta(nw)$.

y-fast aplica indireccion: solo los representantes de buckets estan en un
x-fast. Cada bucket tiene tamano esperado/amortizado $\Theta(w)$ y se guarda en
una estructura balanceada local. Se consulta el representante en $\bigO(\log w)$
y luego se busca dentro del bucket en $\bigO(\log w)$. El espacio total baja a
$\bigO(n)$.

\subsection*{I4. PS8 MIT: compresion de signatures}
Referencia: \texttt{../originales/psets/ps8.pdf}. Se tienen $k=\log^\epsilon n$
hashes de $\log n$ bits separados por ceros dentro de una palabra. Se quiere
empaquetarlos juntos a la derecha en tiempo $\bigO(1)$ word RAM.

\paragraph{Solucion propia.}
La herramienta conceptual es una multiplicacion por una mascara constante
precomputada que mueve cada campo a su posicion objetivo, seguida de mascaras
para eliminar basura. Como los campos no se solapan despues del desplazamiento,
la suma/multiplicacion paralela dentro de palabra actua como una operacion SIMD:
cada hash $h_i$ se desplaza por una cantidad fija hasta quedar contiguo al
siguiente. La prueba consiste en fijar una mascara con 1s en las posiciones de
inicio de cada campo, multiplicar por un patron que copia esos campos a offsets
controlados, y aplicar una mascara final de $k\log n$ bits. Todas las operaciones
son de palabra, por lo que el tiempo es $\bigO(1)$.

\section{Static Trees, RMQ y LCA}

\subsection*{S1. RMQ a LCA con Cartesian tree}
Dado un arreglo $A[1..n]$, construye su Cartesian tree min-heap. Prueba que
$\RMQ_A(i,j)$ es el LCA de los nodos $i$ y $j$.

\paragraph{Solucion.}
El Cartesian tree mantiene inorder igual al orden de indices y heap property por
valor. Para cualquier intervalo $[i,j]$, el nodo de minimo valor en ese intervalo
es el primer nodo del intervalo que separa el inorder en izquierda/derecha; por
tanto es ancestro de todos los indices del intervalo. En particular, es el
ancestro comun mas bajo de $i$ y $j$. Si hubiera un ancestro comun mas bajo con
menor profundidad dentro del intervalo, violaria la propiedad de heap o la
separacion inorder. Entonces $\RMQ_A(i,j)=\LCA_T(i,j)$.

\subsection*{S2. LCA a +-1 RMQ}
Explica la reduccion de LCA en un arbol estatico a RMQ sobre un arreglo con
diferencias consecutivas $\pm 1$.

\paragraph{Solucion.}
Hacer Euler tour del arbol y guardar dos arreglos: $E$, nodos visitados, y $D$,
profundidades. Tambien guardar $first[v]$, primera aparicion de $v$ en $E$. Para
una query $\LCA(u,v)$, sea $l=\min(first[u],first[v])$ y $r=\max(...)$. El LCA es
el nodo $E[t]$ donde $D[t]$ es minimo en $[l,r]$. Como el Euler tour se mueve por
una arista en cada paso, las profundidades cambian exactamente en $+1$ o $-1$.
Asi el problema se reduce a $\pm 1$ RMQ.

\subsection*{S3. PS10 MIT: succinct +-1 RMQ}
Referencia: \texttt{../originales/psets/ps10.pdf}. Disena una estructura
succinct para $\pm 1$ RMQ con $n+o(n)$ bits y query constante.

\paragraph{Solucion propia.}
Usar dos niveles. Primero codificar el arreglo $\pm 1$ por sus diferencias:
un bit por paso, total $n+o(n)$ bits. Dividir en bloques pequenos de tamano
$b=\frac12\log n$. Cada bloque tiene solo $2^{b-1}=\sqrt n$ formas posibles; se
precomputan tablas de respuestas internas para cada forma usando
$o(n)$ bits. Luego agrupar bloques en super-bloques y guardar minimos de bloque
en una macroestructura RMQ sparse/succinct sobre los minimos. Una query se parte
en: sufijo del bloque izquierdo, bloques completos intermedios, prefijo del
bloque derecho. Los extremos se responden por microtabla; el centro por macro
RMQ. El espacio extra es $o(n)$ y el tiempo es $\bigO(1)$.

\subsection*{S4. Imperial Roads}
Se tiene un grafo ponderado. Para cada carretera consultada, se pide el peso del
MST si esa carretera debe estar incluida.

\paragraph{Solucion.}
Construir un MST $T$ con Kruskal y peso total $W$. Preprocesar $T$ para responder
$maxEdge(u,v)$ en el camino entre dos nodos, usando binary lifting, HLD o LCA con
RMQ. Para una arista $e=(u,v,w)$, al agregarla a $T$ se forma un ciclo. Para
mantener un arbol con $e$, se debe eliminar una arista del camino $u\leadsto v$.
La mejor eleccion es eliminar la de mayor peso. Por tanto
\[
ans(e)=W+w-maxEdge(u,v).
\]
Correctitud: por la propiedad del ciclo del MST, en cualquier ciclo, una arista
de peso maximo puede excluirse de algun MST. Complejidad: Kruskal
$\bigO(m\log m)$, preproceso $\bigO(n\log n)$, query $\bigO(\log n)$ con binary
lifting o $\bigO(1)$ si se reduce a RMQ apropiadamente.

\section{Strings}

\subsection*{T1. Buscar un patron con suffix array}
Dado texto fijo $S$ y patron $P$, explica como reportar ocurrencias usando
suffix array.

\paragraph{Solucion.}
Construir $\SA$ de $S\$$. Los sufijos que empiezan con $P$ forman un intervalo
contiguo en orden lexicografico, porque todos tienen prefijo comun $P$ y ningun
sufijo fuera del rango puede quedar entre ellos sin compartir el mismo prefijo.
Hacer dos binary searches: primera posicion con sufijo $\ge P$ y primera posicion
con sufijo $>P$ en orden prefijo. El intervalo resultante contiene las posiciones
de ocurrencia. Tiempo de construccion segun algoritmo: $\bigO(n\log n)$ con
Manber-Myers o $\bigO(n)$ con DC3/SA-IS; query directa
$\bigO(|P|\log n + occ)$.

\subsection*{T2. Contar subcadenas distintas}
Prueba la formula usando suffix array y LCP.

\paragraph{Solucion.}
El sufijo que empieza en $i$ aporta $n-i$ prefijos, todos subcadenas de $S$.
Si procesamos sufijos en orden de $\SA$, el sufijo actual comparte exactamente
$\lcp$ subcadenas-prefijo con el sufijo anterior; esas ya fueron contadas. Por
tanto el numero de nuevas subcadenas aportadas es
$(n-\SA[i])-\lcp[i]$. Sumando:
\[
\#distinct = \frac{n(n+1)}2 - \sum_i \lcp[i].
\]

\subsection*{T3. Longest common substring}
Dados strings $A$ y $B$, encuentra la subcadena comun mas larga.

\paragraph{Solucion.}
Construir $S=A\#B\$$ con separadores que no aparecen en los strings. Construir
$\SA$ y $LCP$. La respuesta es el maximo $LCP[i]$ tal que los sufijos
$\SA[i-1]$ y $\SA[i]$ vienen de strings distintos. Si dos sufijos tienen un
prefijo comun de longitud $L$, ese prefijo es una subcadena comun; el maximo
entre pares adyacentes basta porque, en un conjunto de sufijos con prefijo comun,
aparecen consecutivos en $\SA$.

\subsection*{T4. Manber-Myers}
Explica el algoritmo de doubling para construir suffix array.

\paragraph{Solucion.}
Mantener clases $rank_k[i]$ para subcadenas de longitud $2^k$ que empiezan en
$i$. Inicialmente se ordena por caracter. Para pasar de $k$ a $k+1$, ordenar cada
posicion $i$ por el par
\[
(rank_k[i], rank_k[i+2^k]).
\]
Despues de ordenar, se reasignan clases. Cada iteracion duplica la longitud
comparada; tras $\lceil\log n\rceil$ iteraciones, las clases identifican sufijos
completos. Con counting/radix sort sobre ranks, cada iteracion cuesta
$\bigO(n)$ y el total es $\bigO(n\log n)$.

\subsection*{T5. DC3 en una pagina}
Da el esquema de DC3 y su recurrencia.

\paragraph{Solucion.}
Particionar posiciones por residuo modulo 3: $S_0,S_1,S_2$. Ordenar
recursivamente sufijos en $S_1\cup S_2$ representando cada posicion por ternas
de caracteres y convirtiendolas a ranks. Con esos ranks, ordenar $S_0$ por
$(S[i], rank[i+1])$. Luego hacer merge lineal entre $S_0$ y $S_1\cup S_2$:
para comparar $i\in S_0$ con $j\in S_1$, usar
$(S[i],rank[i+1])$; para $j\in S_2$, usar
$(S[i],S[i+1],rank[i+2])$. La recurrencia es
$T(n)=T(2n/3)+\bigO(n)=\bigO(n)$.

\section{Dynamic Graph}

\subsection*{D1. HLD path query}
Prueba que HLD produce solo $\bigO(\log n)$ cadenas en cualquier camino
raiz-nodo y da path query.

\paragraph{Solucion.}
Una arista ligera hacia un hijo $v$ cumple $size(v)\le size(u)/2$. Cada vez que
se baja por una arista ligera, el tamano del subarbol se reduce al menos a la
mitad, asi que no puede ocurrir mas de $\log n$ veces en una ruta. Las aristas
pesadas forman cadenas. Para query path $u,v$, mientras las cabezas de cadena
sean distintas, subir la cabeza mas profunda y consultar el segmento
$[in[head],in[node]]$. Al final ambos estan en la misma cadena y se consulta el
segmento restante. Con segment tree, cada segmento cuesta $\bigO(\log n)$ y hay
$\bigO(\log n)$ segmentos: total $\bigO(\log^2 n)$.

\subsection*{D2. PS11 MIT: LCA con link-cut tree}
Referencia: \texttt{../originales/psets/ps11.pdf}. Explica como responder LCA
en $\bigO(\log n)$ usando link-cut trees.

\paragraph{Solucion propia.}
En un link-cut tree, \texttt{access(v)} expone la ruta desde la raiz representada
hasta $v$ como preferred path. Si primero hacemos \texttt{access(v)}, la ruta
raiz-$v$ queda marcada como la ultima ruta accedida. Luego hacemos
\texttt{access(w)}. Durante este segundo acceso, el punto donde la busqueda
cambia de una preferred path ya existente a la ruta hacia $w$ identifica el
ultimo vertice comun entre raiz-$v$ y raiz-$w$, es decir, $\LCA(v,w)$. Muchas
implementaciones devuelven precisamente el ultimo nodo tocado/reconectado por
\texttt{access(w)} despues de \texttt{access(v)}. Como cada access cuesta
$\bigO(\log n)$ amortizado, la query LCA tambien.

\subsection*{D3. Euler-tour tree para component sum}
Mantener un forest dinamico con pesos en vertices: \texttt{link}, \texttt{cut},
\texttt{addVertex(v,x)} y \texttt{componentSum(v)}.

\paragraph{Solucion.}
Representar cada arbol por su Euler tour en un treap/BST implicito. Cada
ocurrencia de vertice apunta al nodo del treap; guardar en cada treap la suma
agregada de pesos. Para no contar un vertice multiples veces, se puede asignar
el peso real solo a una ocurrencia canonica y peso cero a las demas. \texttt{link}
y \texttt{cut} son secuencias de splits/merges en $\bigO(\log n)$. Para
\texttt{componentSum(v)}, tomar cualquier ocurrencia de $v$, encontrar la raiz
del treap y leer su suma agregada. \texttt{addVertex} actualiza la ocurrencia
canonica y recalcula agregados en $\bigO(\log n)$.

\subsection*{D4. Fully dynamic connectivity general}
Bosqueja la estructura con niveles y Euler-tour trees.

\paragraph{Solucion.}
Mantener niveles $0,\ldots,\lfloor\log n\rfloor$. En cada nivel hay un spanning
forest de las aristas de ese nivel o superiores, representado con Euler-tour
trees. Las aristas tree mantienen conectividad; las non-tree son candidatas de
reemplazo. Al borrar una arista que no es tree, solo se elimina. Al borrar una
tree edge, el arbol se parte en dos componentes; se buscan aristas non-tree que
reconecten ambas partes, promoviendo/demoviendo aristas por niveles para amortizar
el costo. La intuicion del bound es que una arista solo puede bajar/subir
$\bigO(\log n)$ niveles y cada operacion de ETT cuesta $\bigO(\log n)$, dando
$\bigO(\log^2 n)$ amortizado para updates.

\section*{Tabla final de complejidades}
\begin{center}
\begin{tabular}{lll}
\toprule
Tema & Estructura & Complejidad clave \\
\midrule
Integer & vEB & $\bigO(\log w)$ query/update, espacio grande \\
Integer & y-fast & $\bigO(\log w)$ amortizado, espacio $\bigO(n)$ \\
Integer & fusion tree & $\bigO(\log_w n)$ predecessor \\
Static trees & LCA via RMQ & $\bigO(n)$ preprocess, $\bigO(1)$ query \\
Strings & Manber-Myers & $\bigO(n\log n)$ construccion SA \\
Strings & DC3/SA-IS & $\bigO(n)$ construccion SA \\
Dynamic graph & HLD & $\bigO(\log^2 n)$ path query/update \\
Dynamic graph & link-cut & $\bigO(\log n)$ amortizado \\
Dynamic graph & Euler-tour tree & $\bigO(\log n)$ link/cut/connectivity forest \\
\bottomrule
\end{tabular}
\end{center}

\end{document}
